\section{主要研究工作总结}
\label{sec:ch07_research_work_summary}

以下按第\ref{sec:research_contents}节的四项具体工作总结，其中抓取示教与数据质量控制共同对应专家示教与数据主线。针对灵巧手操作中机器人运动、接触感知与任务执行需要协同组织的问题，本文构建了RM75B机械臂与途见科技自研欠驱动灵巧手的多任务仿真平台。在MuJoCo环境中集成机器人模型、控制器、任务逻辑和触觉记录模块，结合曲面几何布局建立16个Patch、712维processed触觉观测体系，明确感知区域、原生读取与信号处理之间的关系。通过多任务接口及指定历史案例的观测检查，形成支撑后续抓取和触觉研究的仿真基础。这部分工作侧重系统集成与感知建模，统一了后续抓取和触觉研究所用的执行接口与观测定义。

针对几何抓取候选不能直接转化为机器人可执行操作的问题，本文适配GraspQP官方抓取评价指标，参考DexEvolve设计实现进化搜索，并结合腕部姿态扩展形成多来源候选。通过机械臂逆运动学和连续轨迹求解连接分阶段专家示教，并按第\ref{sec:grasp_reachability}节所述，结合桌面几何惩罚、局部接近约束、执行监测和尾段Strict检查验证轨迹与任务结果，区分操作完成与示教质量。局部扰动轨迹与失败分类检验了该流程在已有案例中的表现，所得专家轨迹也是后续数据记录的来源。工作重点在于将已有候选方法与机器人约束、任务执行和示教质量检查结合起来。

在专家示教基础上，本文围绕数据语义、采集时序和来源追踪建立多源数据组织与质量控制体系。通过区分策略观测与完整仿真状态、位置命令与逆运动学目标，明确状态和动作合同，检查动作命令与执行前后观测的时序配对关系，并完成NPZ、Manifest与LeRobot之间的转换组织。质量检查覆盖数值有效性、连续性、完整性、版本和重复关系，既核验了指定修复数据的字段映射，也保留旧记录中的命令不连续与状态字段缺失问题。这些检查分别对应采集、转换和训练准入中的不同问题，使指定样本的修复结果、历史缺陷及来源关系能够按版本追踪。

针对仿真与真实触觉存在分布差异且缺少认证同步物理配对的问题，本文利用真实统计先验与仿真探针构建伪配对，分别学习独立双域VAE表示，再由GRU Translator建立潜变量映射。在此基础上研究共享参数向量场、对抗优化及多视角判别与残差修正，并结合分布距离、邻域覆盖、输入依赖和判别器跨Session可靠性开展综合评价。这些工作使用单独登记并冻结的触觉子集，沿用前述数据管理与质量规范；LeRobot示教导出则保持其独立的数据用途。研究重点由构造统计监督，逐步转向检验重建结构、输入依赖和评价器的可信度。现有实验获得了限定指标下的局部改善，同时揭示接受率、分布质量和评价可靠性的差异，尚未完成完整触觉Sim-to-Real迁移验收。
